\begin{table}
\caption{Clinical Diagnostic Strengths, Limitations and Usability Guidelines for Multi-Modal XAI}
\label{tab:table_5_03_02_clinical_diagnostic_strengths_and_limitations_of_xai_methods}
\begin{tabular}{llllll}
\toprule
XAI Technique & Optimal Clinical Scenario & Radiologist Usability Rating & Diagnostic Action Triggered & Failure Mode / Hazard & Remediation Protocol \\
\midrule
Grad-CAM / Grad-CAM++ & Screening Mammography & Breast Ultrasound & 9.4 / 10 & Validates BIRADS mass location and verifies margins & Can highlight benign glandular tissue surrounding small spiculations & Pair with Integrated Gradients to check fine structural details \\
Integrated Gradients & High-Resolution Histopathology Biopsy Patches & 8.8 / 10 & Identifies nuclear atypia, mitoses, and tubule loss & High-frequency background artifact can confuse non-expert observers & Apply SmoothGrad Gaussian blur filtering (sigma=1.0) \\
Kernel SHAP & Clinical Demographics & Longitudinal EHR Panels & 9.6 / 10 & Directly explains patient risk factors (Age, Parity, Hormone Therapy) & Correlated clinical features may split attribution mass & Group collinear laboratory markers into hierarchical clusters \\
Cross-Attention Rollout & Multidisciplinary Tumor Board (MDT) Review & 9.5 / 10 & Explains which modality broke diagnostic deadlocks in discordant cases & Abstract modality graph requires brief clinician orientation & Render intuitive chord diagrams with normalized percentage weights \\
\bottomrule
\end{tabular}
\end{table}
